%!TEX program = uplatex
\RequirePackage{plautopatch}
\documentclass[uplatex,dvipdfmx]{jlreq}
\usepackage{amsmath,amssymb}

\begin{document}
\thispagestyle{empty}

\noindent
{\large\bfseries 計算図表、Web、微分方程式の標準化とツイスター理論}

\medskip
\hspace*{\fill}{\large\bfseries 佐藤 肇（名古屋大学・多元数理）}

\bigskip\bigskip

最近のデジタル型コンピューターの発展により影がうすくなったが、
もう一度アナログ型コンピューターを考えてみよう。
もう中年にしか記憶にないかもしれないが、計算尺を知っている人も多いであろう。
計算図表は計算尺をやや複雑にしたものであり、Web 理論の基となった。

\bigskip

\noindent\textbf{問題：}$(x,y,z)$ 空間での関数 $F(x,y,z)$ が与えられた時、
$F(x,y,z) = 0$ を満たす組 $(x,y,z)$ を求めよ。
あるいは $(x_0, y_0)$ に対して $F(x_0, y_0, z_0) = 0$ を満たす $z_0$ を探せといってもよい。

\bigskip

\noindent\textbf{例１}\quad $F(x,y,z) = z^2 + yz + x$

この時は、積が $x$、和が $-y$ となる２つの数 $z_1, z_2$ を求めよという問題となる。

\medskip

\noindent\textbf{例２}\quad $F(x,y,z) = x^2 + y^2 + z$

\bigskip

\noindent\textbf{解法１}\quad $F(x,y,\varphi(x,y)) = 0$ となる陰関数 $\varphi(x,y)$ を求めればよい。

\medskip

より図形的にこのような解を求める方法は次のようになるであろう。

\medskip

\noindent\textbf{解法２}\quad $F(x,y,z) = 0$ は空間の一つの曲面を定める。
その曲面の $(x_0, y_0)$ における高さが求める $z_0$ の値となる。

\medskip

\noindent\textbf{解法３}\quad さらに、$(x,y)$ 平面上で解を知るには、
曲面 $F(x,y,z) = 0$ の $z\,(= \varphi(x,y)) = \mathrm{const}$ という等高線を細かく引いて、
$x = x_0$、$y = y_0$ という直線と交わる等高線の示している高さ $z_0$ をみればよい。
この方法は、$F(x,y,z) = 0$ という曲面を与えると、
\begin{itemize}
  \setlength{\itemsep}{2pt}
  \item $x = \mathrm{const}$ という直線による一組の族、
  \item $y = \mathrm{const}$ という直線による一組の族、
  \item $z = \mathrm{const}$ という曲線による一組の等高線の族
\end{itemize}
が与えられているとみなしている。

これが共点図表で、いいかえれば、平面の $3$-Web でもある。

\bigskip

例１の場合を考えると、$z = \mathrm{const}$ という曲線もすべて直線になり、
３組の直線族による $3$-Web となる。
例２はそうではないが、
平面の $x \mapsto x^2$、$y \mapsto y^2$ という座標変換により直線族となる。

\medskip

このように、曲線の族による平面の $3$-Web を座標変換により直線族による $3$-Web と
みなすことができるかというのが、
共線図表（＝計算図表（狭い意味で））化可能性という計算図表学の基本問題であり、
また Web 理論・微分方程式論の中心問題でもある。
この解答は、ツイスター理論と密接な関係があることがわかる。

\end{document}